\chapter{Przetwarzanie danych strumieniowych}
\label{ch:przetwarzanie_danych}

Niniejszy rozdział przedstawia teoretyczne podstawy przetwarzania danych strumieniowych, ze~szczególnym uwzględnieniem
specyfiki rynków finansowych i~kryptowalutowych.
Omówiono kluczowe koncepcje, modele przetwarzania oraz wyzwania związane z~obsługą ciągłych przepływów danych
w~czasie rzeczywistym.


\section{Charakterystyka danych strumieniowych}
\label{sec:charakterystyka_strumieni}

Dane strumieniowe (\textit{streaming data}) to~ciągły, potencjalnie nieskończony przepływ rekordów generowanych
w~czasie rzeczywistym przez różnorodne źródła.
W~odróżnieniu od~tradycyjnych zbiorów danych o~skończonym rozmiarze, strumienie charakteryzują się brakiem określonego
początku i~końca --- nowe dane napływają nieprzerwanie, a~ich przetwarzanie musi odbywać się na~bieżąco~\cite{data_stream_def}.

W~kontekście rynków kryptowalutowych dane strumieniowe obejmują przede wszystkim zdarzenia transakcyjne
(\textit{trade events}), aktualizacje arkuszy zleceń (\textit{order book updates}) oraz zmiany cen i~wolumenów.
Giełdy takie jak Binance generują tysiące komunikatów na~sekundę dla~popularnych par handlowych, co~stawia wysokie
wymagania przed systemami przetwarzającymi.

Kluczowe cechy danych strumieniowych z~rynków finansowych to~\cite{burst_events, time_value_curve, 8_requirements}:
\begin{itemize}
    \item \textbf{Wysoka częstotliwość} --- komunikaty napływają w~odstępach milisekundowych, szczególnie w~okresach
    zwiększonej zmienności rynkowej.
    \item \textbf{Zmienność natężenia} --- przepływ danych nie~jest stały; w~momentach publikacji istotnych informacji
    rynkowych wolumen komunikatów może wzrosnąć wielokrotnie.
    \item \textbf{Krytyczność czasowa} --- wartość informacji maleje wraz z~upływem czasu; opóźnione dane tracą
    użyteczność dla~algorytmów handlowych.
\end{itemize}


\section{Modele przetwarzania danych}
\label{sec:modele_przetwarzania}

W~literaturze i~praktyce inżynierskiej wyróżnia się dwa podstawowe podejścia do~przetwarzania danych: przetwarzanie
wsadowe (\textit{batch processing}) oraz przetwarzanie strumieniowe (\textit{stream processing}).
Każde z~nich charakteryzuje się odmiennymi właściwościami i~znajduje zastosowanie w~różnych scenariuszach.

\subsection{Przetwarzanie wsadowe}
\label{subsec:batch_processing}

Przetwarzanie wsadowe polega na~gromadzeniu danych przez~określony czas, a~następnie ich jednorazowym przetworzeniu
jako~skończonego zbioru.
Model ten jest dobrze ugruntowany i~szeroko stosowany w~systemach analitycznych, gdzie~akceptowalne są~opóźnienia
rzędu minut, godzin lub~dni.

Typowy przepływ przetwarzania wsadowego obejmuje ekstrakcję danych ze~źródeł, ich~transformację zgodnie z~logiką
biznesową oraz~załadowanie do~docelowego magazynu danych (\textit{ETL --- Extract, Transform, Load}).
Procesy te~uruchamiane są~cyklicznie, według ustalonego harmonogramu.

Zalety przetwarzania wsadowego to~prostota implementacji, łatwość debugowania oraz~możliwość ponownego przetworzenia
całego zbioru w~przypadku wykrycia błędów.
Wadą jest natomiast opóźnienie między~powstaniem danych a~ich dostępnością dla~odbiorców, co~w~kontekście rynków
finansowych może być nieakceptowalne~\cite{bigdata_principles}.

\subsection{Przetwarzanie strumieniowe}
\label{subsec:stream_processing}

Przetwarzanie strumieniowe zakłada przetwarzanie danych natychmiast po~ich nadejściu, bez~oczekiwania na~zgromadzenie
pełnego zbioru.
Każdy komunikat jest analizowany indywidualnie lub~w~kontekście niewielkiego okna czasowego, a~wyniki są~udostępniane
z~minimalnym opóźnieniem.

W~tym modelu dane przepływają przez~zdefiniowany potok przetwarzania (\textit{pipeline}), składający się z~operatorów
realizujących transformacje, filtrowanie, agregację czy~wzbogacanie rekordów.
Strumień wyjściowy może być kierowany do~kolejnych etapów przetwarzania, zapisywany w~bazie danych lub~wykorzystywany
do~generowania alertów.

Przetwarzanie strumieniowe jest szczególnie istotne w~kontekście danych finansowych, gdzie~opóźnienia rzędu
milisekund mogą decydować o~rentowności strategii handlowych.
Wymaga jednak bardziej złożonej infrastruktury oraz~starannego zarządzania stanem przetwarzania~\cite{8_requirements}.

\subsection[Architektura Lambda i Kappa]{Architektura Lambda i~Kappa}
\label{subsec:lambda_kappa}

W~praktyce często stosuje się architektury hybrydowe, łączące zalety obu podejść.
Architektura Lambda zakłada równoległe działanie dwóch ścieżek przetwarzania: warstwy wsadowej (\textit{batch layer})
zapewniającej kompletność i~poprawność danych historycznych oraz~warstwy strumieniowej (\textit{speed layer})
dostarczającej wyniki w~czasie rzeczywistym.
Warstwa udostępniania (\textit{serving layer}) łączy wyniki z~obu ścieżek, udostępniając spójny widok danych~\cite{bigdata_principles}.

Architektura Kappa stanowi uproszczenie tego modelu poprzez eliminację dedykowanej warstwy wsadowej.
W~tym podejściu wszystkie dane są~przetwarzane jako~strumień, a~przetwarzanie historyczne realizowane jest poprzez
ponowne odtworzenie strumienia z~trwałego magazynu zdarzeń.
Model ten jest koncepcyjnie prostszy, lecz~wymaga odpowiedniej infrastruktury do~przechowywania i~odtwarzania
historycznych strumieni~\cite{kappa_architecture}.


\section{Semantyka dostarczania komunikatów}
\label{sec:semantyka_dostarczania}

Kluczowym aspektem systemów przetwarzania strumieniowego jest semantyka dostarczania komunikatów, określająca
gwarancje dotyczące tego, ile~razy każdy komunikat zostanie przetworzony.
Wyróżnia się trzy podstawowe poziomy gwarancji:

\subsection[At-most-once (co najwyżej raz)]{At-most-once (co~najwyżej raz)}
\label{subsec:at_most_once}

W~semantyce \textit{at-most-once} komunikat może zostać utracony, lecz~nigdy nie~zostanie przetworzony więcej
niż~jeden raz.
Jest to~najprostsza do~implementacji gwarancja, gdyż~nie wymaga mechanizmów potwierdzeń ani~powtórzeń.
Stosowana jest w~scenariuszach, gdzie~sporadyczna utrata danych jest~akceptowalna, na~przykład w~systemach
monitoringu, gdzie~pojedyncze brakujące pomiary nie~wpływają istotnie na~ogólny obraz.

W~kontekście danych finansowych semantyka ta~jest rzadko stosowana, ponieważ utrata nawet pojedynczej transakcji
może prowadzić do~niespójności w~analizach lub~błędnych decyzji inwestycyjnych~\cite{at_least_once}.

\subsection[At-least-once (co najmniej raz)]{At-least-once (co~najmniej raz)}
\label{subsec:at_least_once}

Semantyka \textit{at-least-once} gwarantuje, że~każdy komunikat zostanie przetworzony przynajmniej jeden raz,
lecz~dopuszcza możliwość duplikatów.
Realizowana jest poprzez mechanizm potwierdzeń (\textit{acknowledgements}) --- komunikat pozostaje w~kolejce do~czasu
otrzymania potwierdzenia pomyślnego przetworzenia.
W~przypadku awarii konsumenta lub~braku potwierdzenia, komunikat jest dostarczany ponownie.

Model ten wymaga, aby~konsumenci byli idempotentni, czyli~wielokrotne przetworzenie tego samego komunikatu
nie~powinno prowadzić do~niepożądanych efektów ubocznych.
W~praktyce idempotentność osiąga się poprzez deduplikację na~podstawie unikalnych identyfikatorów komunikatów
lub~projektowanie operacji w~sposób naturalnie idempotentny (np.~ustawianie wartości zamiast inkrementacji)~\cite{at_least_once}.

\subsection{Exactly-once (dokładnie raz)}
\label{subsec:exactly_once}

Semantyka \textit{exactly-once} stanowi najsilniejszą gwarancję --- każdy komunikat jest przetwarzany dokładnie
jeden raz, bez~utraty i~bez~duplikatów.
Jest to~najbardziej pożądana, lecz~zarazem najtrudniejsza do~osiągnięcia semantyka.

Realizacja wymaga koordynacji między~producentem, brokerem i~konsumentem, często z~wykorzystaniem mechanizmów
transakcyjnych lub~rozproszonego zatwierdzania (\textit{distributed commit}).
Ze~względu na~złożoność implementacji oraz~wpływ na~wydajność, w~wielu systemach stosuje się semantykę
\textit{at-least-once} w~połączeniu z~idempotentnym przetwarzaniem jako~praktyczny kompromis~\cite{at_least_once}.


\section[Okna czasowe w przetwarzaniu strumieniowym]{Okna czasowe w~przetwarzaniu strumieniowym}
\label{sec:okna_czasowe}

Wiele operacji analitycznych wymaga agregacji danych w~określonych przedziałach czasowych.
W~przetwarzaniu strumieniowym wykorzystuje się koncepcję okien czasowych, które~grupują zdarzenia na~potrzeby obliczeń zbiorczych.

\subsection{Okna stałe (fixed windows)}
\label{subsec:fixed_windows}

Okna stałe dzielą strumień na~rozłączne, następujące po~sobie przedziały o~jednakowej długości.
Każde zdarzenie należy do~dokładnie jednego okna, a~wyniki agregacji są~emitowane po~zamknięciu każdego okna.

Przykładem zastosowania jest obliczanie średniej ceny w~interwałach minutowych --- każda minuta stanowi osobne okno,
a~po jej~zakończeniu publikowany jest wynik (np.~cena otwarcia, zamknięcia, maksimum, minimum).
Takie agregaty stanowią podstawę wykresów świecowych (\textit{candlestick charts}) powszechnie stosowanych
w~analizie technicznej~\cite{dataflow_model}.

\subsection{Okna przesuwne (sliding windows)}
\label{subsec:sliding_windows}

Okna przesuwne obejmują zdarzenia z~określonego przedziału czasowego, lecz~w~odróżnieniu od~okien stałych,
nakładają się na~siebie.
Okno przesuwa się wraz z~nadejściem nowych zdarzeń, umożliwiając ciągłą aktualizację wyników agregacji.

Typowym zastosowaniem jest obliczanie średniej kroczącej (\textit{moving average}), gdzie~wynik jest aktualizowany
przy~każdym nowym zdarzeniu, uwzględniając ostatnie N~minut lub~ostatnie N~transakcji.
Okna przesuwne zapewniają płynniejszy obraz zmian, lecz~wymagają większych zasobów obliczeniowych ze~względu
na~nakładanie się przedziałów~\cite{dataflow_model}.

\subsection{Okna sesyjne (session windows)}
\label{subsec:session_windows}

Okna sesyjne grupują zdarzenia na~podstawie aktywności, a~nie sztywnych granic czasowych.
Sesja rozpoczyna się wraz z~pierwszym zdarzeniem i~trwa do~momentu, gdy~przez~zdefiniowany okres nie~napłyną żadne
nowe zdarzenia (tzw.~\textit{gap}).
Po~przekroczeniu tego progu sesja jest zamykana, a~kolejne zdarzenie rozpoczyna nową sesję.

W~kontekście rynków finansowych okna sesyjne mogą być wykorzystywane do~identyfikacji okresów wzmożonej aktywności
handlowej lub~grupowania serii szybkich transakcji wykonywanych przez~algorytmy wysokiej częstotliwości~\cite{dataflow_model}.


\section[Problemy i wyzwania]{Problemy i~wyzwania}
\label{sec:problemy_strumieniowe}

Przetwarzanie danych strumieniowych wiąże się z~szeregiem wyzwań technicznych i~koncepcyjnych, które~wymagają
starannego rozwiązania na~etapie projektowania systemu.

\subsection{Zarządzanie czasem zdarzeń}
\label{subsec:czas_zdarzen}

W~systemach rozproszonych rozróżnia się czas zdarzenia (\textit{event time}), czyli~moment faktycznego wystąpienia
zdarzenia u~źródła, oraz~czas przetwarzania (\textit{processing time}), czyli~moment dotarcia zdarzenia do~systemu.
Różnica między~tymi znacznikami wynika z~opóźnień sieciowych, kolejkowania oraz~ewentualnych awarii.

Zdarzenia mogą docierać w~kolejności innej niż~chronologiczna (tzw.~\textit{out-of-order events}), co~komplikuje
obliczenia oparte na~czasie.
System musi decydować, jak~długo czekać na~potencjalnie opóźnione zdarzenia przed zamknięciem okna czasowego
i~emisją wyników.
Zbyt~krótkie oczekiwanie prowadzi do~pominięcia opóźnionych danych, zbyt~długie --- do~zwiększenia opóźnienia wyników.

Mechanizm \textit{watermarks} pozwala śledzić postęp czasu zdarzeń i~sygnalizować, że~wszystkie zdarzenia
o~znaczniku czasowym mniejszym od~pewnej wartości zostały już~odebrane.
Zdarzenia przychodzące po~przekroczeniu znacznika czasu są~traktowane jako~spóźnione i~mogą być odrzucane
lub~przetwarzane w~specjalny sposób~\cite{dataflow_model}.

\subsection{Zarządzanie stanem}
\label{subsec:zarzadzanie_stanem}

Wiele operacji strumieniowych wymaga utrzymywania stanu między~kolejnymi komunikatami --- na~przykład obliczenie
średniej kroczącej wymaga pamiętania poprzednich wartości, a~agregacja w~oknie czasowym --- gromadzenia zdarzeń
do~momentu zamknięcia okna.

Zarządzanie stanem w~systemach rozproszonych jest wyzwaniem ze~względu na~konieczność zapewnienia trwałości
(stan nie~może zostać utracony w~przypadku awarii), spójności (wszystkie repliki muszą widzieć ten~sam stan)
oraz~wydajności (dostęp do~stanu nie~może stanowić wąskiego gardła).

Stosowane podejścia obejmują lokalne magazyny stanu z~mechanizmem punktów kontrolnych (\textit{checkpointing}),
zewnętrzne bazy danych o~niskich opóźnieniach oraz~rozproszone struktury danych z~replikacją~\cite{flink}.

\subsection[Backpressure i kontrola przepływu]{Backpressure i~kontrola przepływu}
\label{subsec:backpressure}

Sytuacja, w~której producent generuje dane szybciej niż~konsument jest w~stanie je~przetworzyć, nazywana jest
przeciążeniem (\textit{backpressure}).
Bez~odpowiednich mechanizmów kontroli przepływu system może ulec destabilizacji --- kolejki rosną nieograniczenie,
pamięć się wyczerpuje, a~opóźnienia stale się~zwiększają.

Strategie radzenia sobie z~backpressure obejmują buforowanie z~ograniczeniem rozmiaru (nadmiarowe komunikaty
są~odrzucane), adaptacyjne spowalnianie producenta (poprzez sygnalizację zwrotną), oraz~skalowanie horyzontalne
konsumentów (dodawanie kolejnych instancji przetwarzających)~\cite{google_sre_book}.

W~kontekście danych z~rynków kryptowalutowych backpressure może wystąpić w~okresach ekstremalnej zmienności,
gdy~wolumen transakcji wielokrotnie przekracza typowe wartości.
System musi być~zaprojektowany tak, aby~degradował wydajność (zwiększając opóźnienie) zamiast tracić dane lub~ulegać
całkowitej awarii.

\subsection[Skalowalność i partycjonowanie]{Skalowalność i~partycjonowanie}
\label{subsec:skalowalnosc}

Skalowanie systemów strumieniowych wymaga możliwości podziału strumienia na~niezależne partycje, które~mogą być
przetwarzane równolegle.
Kluczem partycjonowania jest zazwyczaj atrybut danych zapewniający równomierny rozkład obciążenia oraz~zachowanie
lokalności danych powiązanych logicznie~\cite{scaling}.

W~przypadku danych z~giełd kryptowalut naturalnym kluczem partycjonowania jest symbol pary handlowej --- wszystkie
zdarzenia dotyczące danej pary (np.~BTCUSDT) trafiają do~tej~samej partycji, co~umożliwia lokalne obliczanie
agregatów bez~koordynacji między~partycjami.

Wyzwaniem jest zapewnienie równomiernego rozkładu obciążenia, gdy~popularność poszczególnych par jest zróżnicowana.
Para BTC/USDT generuje znacznie więcej transakcji niż~mniej popularne pary, co~może prowadzić do~nierównomiernego
obciążenia partycji (\textit{hot spots})~\cite{coingecko}.


\section{Architektura sterowana zdarzeniami}
\label{sec:event_driven}

Architektura sterowana zdarzeniami (\textit{Event-Driven Architecture}, EDA) stanowi paradygmat projektowania
systemów, w~którym przepływ sterowania jest determinowany przez~zdarzenia --- sygnały o~wystąpieniu określonych
stanów lub~akcji w~systemie.

W~modelu tym komponenty systemu komunikują się poprzez publikowanie i~subskrybowanie zdarzeń, bez~bezpośrednich
wywołań synchronicznych.
Producent zdarzenia nie~musi znać odbiorców, a~konsumenci mogą dołączać lub~odłączać się dynamicznie,
co~zapewnia luźne powiązanie (\textit{loose coupling}) między~modułami~\cite{edd_pub_sub}.

Centralna rola w~architekturze EDA przypada brokerowi wiadomości, który~odpowiada za~przyjmowanie zdarzeń
od~producentów, ich~trwałe przechowywanie oraz~dostarczanie do~zainteresowanych konsumentów.
Broker zapewnia buforowanie w~przypadku chwilowej niedostępności odbiorców oraz~może realizować dodatkowe
funkcje, takie jak~filtrowanie, routing czy~transformacja komunikatów~\cite{data_broker}.

Zalety architektury sterowanej zdarzeniami obejmują naturalną asynchroniczność (operacje nie~blokują się wzajemnie),
łatwość dodawania nowych konsumentów (bez~modyfikacji istniejących komponentów), oraz~odporność na~awarie
(zdarzenia są~buforowane do~czasu pomyślnego przetworzenia).
Głównymi wadami są~większa złożoność debugowania (brak liniowego przepływu sterowania) oraz~konieczność zarządzania
ostateczną spójnością (\textit{eventual consistency})~\cite{fundamentals}.


\section{Podsumowanie}
\label{sec:podsumowanie_przetwarzanie}

Przetwarzanie danych strumieniowych z~rynków finansowych wymaga starannego doboru modeli, semantyk i~mechanizmów
odpowiadających specyfice tego~obszaru.
Kluczowe znaczenie ma~minimalizacja opóźnień, zapewnienie niezawodności dostarczania oraz~elastyczna obsługa
zmiennego obciążenia.

Architektura sterowana zdarzeniami w~połączeniu z~semantyką dostarczania \textit{at-least-once} i~idempotentnym
przetwarzaniem stanowi praktyczny kompromis między~złożonością implementacji a~wymaganymi gwarancjami.
Mechanizmy okien czasowych umożliwiają realizację analiz agregacyjnych, natomiast odpowiednie strategie
partycjonowania i~kontroli przepływu zapewniają skalowalność systemu.

W~kolejnych rozdziałach przedstawiono, w~jaki sposób opisane koncepcje zostały zastosowane w~projektowanym systemie
do~przetwarzania danych z~giełdy kryptowalut Binance.

